\begin{definition}
Fix the graph, palette, distinguished edge, exact terminal tree \(T\),
generating coloring, positive-rung ETT, hierarchy base, minimum hierarchy
counterexample \(C\), and its
\noderef{HierarchyMinimumCounterexampleContext}.  Fix also a source
\noderef{HierarchyReplacementPresentation} \(P\) and colors
\(\alpha,\beta\).  Put
\[
 j=\noderef{HierarchyPathNumber}(C.\mathit{hierarchy}),\qquad
 B=\noderef{HierarchyFullTailPrefix}(C,j-2),\qquad
 X=\noderef{HierarchyFullTailPrefix}(C,j-1),
\]
and let \(y_\ell\) be the last displayed vertex of \(T\).  A
\emph{Situation 7.3 disjoint datum} consists first of an actual
original-relative replacement presentation \(P_\sigma\) satisfying
\[
 P_\sigma.\mathit{coloring}
 =\noderef{TwoColorComponentSwap}
   (G,P.\mathit{coloring},\alpha,\beta,y_\ell).
\]
The switched coloring is
\noderef{IsProperEdgeColoring},
\noderef{IsOneEdgeUncolored} at exactly the distinguished edge, and
\noderef{TreeColoringInvariant} on \(X\) relative to the source coloring.
Both \(\alpha\) and \(\beta\) belong to
\(\noderef{TreeMissingColors}(P_\sigma.\mathit{coloring},X)\), and the
same exact ordered suffix is a
\noderef{TashkinovAugmentationSequence} from \(X\) to \(T\) under the
switched coloring.  Moreover, \(\beta\) is missing at \(y_\ell\) and on
\(B\).

Write \(H_\sigma=P_\sigma.\mathit{counterexample.hierarchy}\), let \(q\)
be its last stored index, and set
\[
 A_\sigma=\noderef{HierarchyStarredStage}(H_\sigma,q),\qquad
 D_\sigma=\noderef{HierarchyStoredColorsAt}(H_\sigma,q),\qquad
 \Gamma_\sigma=\bigcup_{\eta\in D_\sigma}
 P_\sigma.\mathit{counterexample.good}.\Gamma^q_\eta.
\]
Let \(T_n^\sigma\) be the terminal ladder tree of the replacement ETT and
put, with \(V(-)=\noderef{TreeSequenceVertexSet}(-)\),
\[
 M_0^\sigma=
 \bigcup_{v\in
 V(P_\sigma.\mathit{replacementBase.core})\setminus V(T_n^\sigma)}
 \noderef{MissingColors}(G,P_\sigma.\mathit{coloring},v).
\]
Define the exact tail-color set
\[
 U_j^\sigma=\{c:\ \exists r\in\operatorname{Fin}|T|,
 |A_\sigma|\le r<|A_\sigma|+j,
 P_\sigma.\mathit{coloring}(T.\mathit{edge}(r))=c\}.
\]
Thus every witness \(r\) selects the individual edge identity
\(T.\mathit{edge}(r)\), including among parallel edges.  The datum records
\(\beta\notin\Gamma_\sigma\) and the full paper condition (9): either
\[
 \beta\notin M_0^\sigma,\quad \beta\notin U_j^\sigma,
 \quad\beta\notin\Gamma_\sigma,\quad\beta\notin D_\sigma,
\]
or
\[
 \beta\in D_\sigma\cap
 \noderef{TreeMissingColors}(P_\sigma.\mathit{coloring},B),\qquad
 (P_\sigma.\mathit{counterexample.good}.\Gamma^q_\beta\cup\{\beta\})
       \cap U_j^\sigma=\varnothing.
\]
Finally the datum contains a color \(\theta\) missing, under the switched
coloring, at the last vertex \(y_j\) of
\(\noderef{HierarchyFullTailPrefix}(C,j)\), and satisfying
\(\theta\notin\Gamma_\sigma\).
\end{definition}
